\documentclass{article}

\usepackage{amsmath, amssymb}
\usepackage{hyperref}

\title{AIDEAZ}
\author{}
\date{}

\begin{document}

\maketitle


\section{Analysis of Rare Events through QGA}

\subsubsection*{Parameter-to-Unitary Mapping}
We use geometry as the main identifier of rare events. Since $SU(n)$ admits a geometric interpretation of unitary evolution, using a higher-dimensional group structure (beyond the Bloch sphere $SU(2)/U(1)$) may be useful. For each chromosome with parameters $\theta \in \mathbb{R}^d$, construct
\[
U(\theta) \in SU(n), \qquad n = 2^q \text{ for } q \text{ qubits}.
\]

\subsubsection*{Embedding into the Lie Algebra}
Embed each unitary into the tangent space of $SU(n)$ via
\[
A(\theta) = \log\!\big( U(\theta) \big),
\]
where $A(\theta) \in \mathfrak{su}(n)$ is skew-Hermitian and satisfies
$\mathrm{Tr}(A(\theta)) = 0$. This linearizes the local geometry, which is required for applying Mahalanobis distances.

\paragraph{Note.}
While reading up on this stuff I came across some of the more intrinsic Lie-group geometry on $SU(n)$—things like the geodesic distance $\|\log(U_1^\dagger U_2)\|$.I'm not fully comfortable with it yet, so I stuck with the tangent-space log map for now since it works naturally with the Mahalanobis setup used.


\subsubsection*{Vectorization}
Flatten each tangent-space matrix:
\[
v(\theta) = \mathrm{vec}\!\big( A(\theta) \big),
\]
yielding a real vector needed for computing means and covariances.

\subsubsection*{Covariance Matrix Estimation}
Given a population $\{v(\theta_i)\}_{i=1}^N$, compute
\[
\mu = \frac{1}{N}\sum_{i=1}^N v(\theta_i), \qquad
\Omega = \frac{1}{N-1}\sum_{i=1}^N (v_i - \mu)(v_i - \mu)^{T}.
\]
Here, $\Omega$ is a covariance matrix that stores the multivariate shape of the population around its mean.

\subsubsection*{Mahalanobis Distance and Density}
Define the metric
\[
D(\theta_a,\theta_b)=\sqrt{(v_a-v_b)^T \Omega^{-1}(v_a-v_b)}.
\]

For each $v(\theta_i)$, define the local density
\[
\rho(\theta_i)
=\left|\left\{ j : D(\theta_i,\theta_j) < \delta \right\}\right|,
\]
where $\delta$ is the local neighbourhood radius used only for density evaluation. Following common practice, $\delta$ is chosen as the median Mahalanobis distance to the $k$-th nearest neighbour across the population.

\subsection*{Rare Event Cluster Determination}

A chromosome belongs to a rare-event cluster if:

\begin{enumerate}
    \item $\rho_{\min} \le \rho(\theta_i) \le \rho_{\text{threshold}}$
    \quad (sparse but not single-point noise),
    \item $f_{LL} \le f(\theta_i) \le f_{UL}$
    \quad (mid-fitness valley region),
    \item its Mahalanobis distance from the population mean satisfies
    \[
    D(\theta_i,\mu) > \epsilon,
    \]
\end{enumerate}

where $\epsilon$ is obtained from a statistical confidence interval (e.g., a $\chi^2$ threshold in Mahalanobis space) under the hypothesis
\[
H_0 : \theta_i \text{ belongs to the main population used for convergence}.
\]

\subsection*{Rare-Event Tracking Across Generations}

Once a rare-event cluster $C_g$ is detected in generation $g$, its centroid and shape are stored as the tuple
\[
A_g = (\mu_g, \Omega_g),
\]
where $\mu_g$ is the mean of the vectorized tangent-space points in the cluster, and $\Omega_g$ is the corresponding covariance matrix. This tuple acts as the signature of the rare-event region at generation $g$.

To determine whether a newly detected rare-event cluster in generation $g$ corresponds to the \emph{same} anomaly as one detected in generation $h$, their centroids are compared using a strict Mahalanobis distance defined with respect to the older cluster's shape:
\[
D_{gh} = \sqrt{(\mu_g - \mu_h)^T \Omega_h^{-1} (\mu_g - \mu_h)}.
\]

A match is accepted only if:
\begin{enumerate}
    \item \textbf{Centroid proximity:}
    \[
    D_{gh} < \tau,
    \]
    meaning the new centroid lies well inside the older ellipsoid.
    \item \textbf{Overlap:} At least say 50\% of points in $C_g$ satisfy
    \[
    \sqrt{(v(\theta_i)-\mu_h)^T \Omega_h^{-1}(v(\theta_i)-\mu_h)} < \tau,
    \]
    ensuring substantial overlap with the older rare region.
    \item \textbf{Fitness consistency:} The mean fitness of $C_g$ differs from that of $C_h$ by no more than say
    10\%.
\end{enumerate}

A rare event is confirmed as a recurring anomaly only if these matching conditions hold for at least some 'm' consecutive generations. This prevents accidental matches caused by random fluctuations and ensures the detection
of only stable rare events.



\section{Fitness function as an observable in QGA}

Using an observable as the fitness function aligns this method with variational quantum algorithms (such as VQE and QAOA), where the goal is to find parameters that minimize or maximize $\langle \psi(\theta)|O|\psi(\theta)\rangle$. However, instead of relying on gradients which are prone to barren plateaus, using observables provide a gradient-free global search mechanism that is better for highly nonlinear, correlated quantum landscapes.
The QGA therefore alternates between (a) purely classical evolutionary updates and (b) quantum evaluations of the observable.

\subsection*{Parameter-to-Circuit Mapping}
Each chromosome $\theta \in \mathbb{R}^d$ defines a parametrized variational quantum circuit
\[
U(\theta)
\]
acting on $q$ qubits. This prepares the quantum state
\[
|\psi(\theta)\rangle = U(\theta)\,|0\rangle^{\otimes q}.
\]

The QGA never manipulates quantum states directly; it only manipulates the classical parameters $\theta$. The circuit is invoked strictly for fitness evaluation.

\subsection*{Observable-Based Fitness}
Given a Hermitian observable $O$ that encodes the task's objective structure, the fitness of a chromosome is defined as
\[
f(\theta)
    = \langle \psi(\theta)\,|\, O \,|\, \psi(\theta) \rangle.
\]

Following are choices for $O$ to be useful in their corresponding scenarios:
\begin{itemize}
    \item Pauli operators (e.g.\ $Z^{\otimes q}$) for alignment-based
    objectives,
    \item task-specific measurement operators,
    \item Hamiltonians $H = \sum_i w_i P_i$ for energy minimization,
    \item projectors $|\psi_{\mathrm{target}}\rangle\langle
          \psi_{\mathrm{target}}|$ for overlap-based tasks.
\end{itemize}



\subsection*{Classical–Quantum Alternation Loop}
The full optimization proceeds in alternating stages:
\begin{enumerate}
    \item \textbf{Classical stage:} A population of parameter vectors
    $\{\theta_i^{(t)}\}$ is updated using standard QGA operators (selection, crossover, mutation, and the amplitude-update rule).     These operations act \emph{only} on the classical parameters.

    \item \textbf{Quantum stage:} Each parameter vector is converted to a quantum state via $U(\theta)$, and its fitness is evaluated using the observable expectation
    \[
    f(\theta) = \langle \psi(\theta)|O|\psi(\theta)\rangle.
    \]

\end{enumerate}

 \section{Potential Applications}

The geometric anomaly-detection framework can be applied wherever optimization algorithms exhibit hidden structure, rare modes, or non-trivial convergence paths. Potential applications include:

\begin{itemize}
    \item \textbf{Diagnostics of variational quantum algorithms:}
    Detecting rare intermediate states that appear during VQE or QAOA optimization, helping understand convergence failures or plateaus.

    \item \textbf{Landscape analysis in evolutionary algorithms:} Revealing secondary basins and rare valleys that standard fitness
    evaluations do not capture.

    \item \textbf{Hamiltonian learning and state identification:} Detecting outlier eigenstates or rare measurement patterns during iterative learning procedures.
\end{itemize}

Replacing the classical fitness function with a quantum observable makes QGA compatible with a wide range of QML and variational tasks. Some representative applications include:

\begin{itemize}
    \item \textbf{Variational quantum eigensolvers (VQE):}
    The observable $O$ can be chosen as a molecular or physical Hamiltonian, allowing QGA to optimize the variational parameters without gradients.

    \item \textbf{State preparation:}
    Using projectors as observables, QGA can search for parameters that prepare a target quantum state.

    \item \textbf{Quantum hardware calibration:}
    Selecting circuit parameters or pulse amplitudes that maximize an experimentally measured observable.

    \item \textbf{QML classification tasks:}
    Encoding class structure in an observable and letting QGA discover parameters that best separate data through quantum measurements.
\end{itemize}

\end{document}
